\section{Introduction and rationale}
\label{sec:introduccion}

Air quality in Nuevo León is a key determinant of public health: daily
exposure to high pollutant concentrations poses a cumulative risk to the
population. The continuous measurements of the Integrated Environmental
Monitoring System (SIMA), described in \cref{sec:sima}, make it possible to
characterize the atmosphere of the Monterrey Metropolitan Area (MMA) and are
the data source of this work. To report air quality, SIMA converts hourly
measurements into a categorical index with a color associated with its
maximum value \parencite{aireNL2026}.

Of the pollutants measured by the network, tropospheric ozone is the one
that motivates this study: it is a secondary pollutant, formed rather than
emitted, and its formation depends jointly on chemical precursors and on
meteorological conditions. That dual dependence is what makes a
multivariate analysis appropriate: no single variable explains an
exceedance day.

\subsection{Tropospheric ozone}
\label{sec:objetivo}

Ozone (\ce{O3}) is a molecule of three oxygen atoms found both in the upper
atmosphere (stratospheric ozone) and at ground level (tropospheric ozone).
Stratospheric ozone is produced naturally by the interaction of molecular
oxygen (\ce{O2}) with ultraviolet radiation, and it attenuates the exposure
of the surface to that radiation \parencite{Bernhard2023-lk}. Tropospheric
ozone, by contrast, is the main ingredient of \textit{smog} and forms through
chemical reactions between nitrogen oxides (\ce{NOx}) and volatile organic
compounds (VOCs) in the presence of heat and sunlight.

Highly industrialized metropolitan areas of Nuevo León have recorded
\ce{O3} and \ce{PM10} levels above the limits of the Mexican standards
NOM-020-SSA1-2021 \parencite{nom020ssa1} and NOM-025-SSA1-2021
\parencite{nom025ssa1}. These trends have been documented for more than
30 years with instruments that measure both ozone (ultraviolet photometry)
and its precursors \ce{NO} and \ce{NO2} (grouped as \ce{NO_x}, by
chemiluminescence), whose photochemical relationship is key to understanding
the formation of the secondary pollutant \parencite{SanfordSillman}.

The MMA is no exception: between 1993 and 2014 tropospheric ozone
concentrations rose steadily, in contrast with the declining trend in Mexico
City over the same period. The increase has been attributed to industrial and
vehicular growth without equivalent control of precursors, and to a
meteorological role modulated by climate change \parencite{atmos15070779}.
The topography of the valley in which the MMA lies aggravates this behavior:
the surrounding mountain ranges favor the stagnation of air masses during
episodes of high pressure, thermal inversion and weak wind, which reduce
dispersion and prolong the exposure of the population.
